\documentclass[11pt]{article}
\usepackage[margin=1in]{geometry}
\usepackage{amsmath,amssymb,amsthm,booktabs,array}
\usepackage[T1]{fontenc}
\usepackage{lmodern}
\usepackage[colorlinks=true,allcolors=blue,breaklinks=true]{hyperref}
\setlength{\emergencystretch}{1.5em}
\newtheorem{theorem}{Theorem}
\newtheorem{lemma}[theorem]{Lemma}
\newtheorem{corollary}[theorem]{Corollary}
\newtheorem{proposition}[theorem]{Proposition}
\newcommand{\NN}{\mathbb N}
\newcommand{\len}{\ell}
\title{Strictness and sibling corrections for OEIS A344611, FORMULA,\protect\\
with an alternative derivation of Manyama's (2026) convolution}
\author{Haobo Ma\\\small ChronoAI Pte Ltd, Singapore
\and Wenlin Zhang\\\small National University of Singapore\\
\small ChronoAI Pte Ltd, Singapore}
\date{November 5, 2026}
\begin{document}
\maketitle
\begin{abstract}
For OEIS A344611, FORMULA (revision 27), we prove
$\mathrm{A344611}(n)\le\allowbreak\mathrm{A160786}(n)$ for every $n\ge0$
and give an alternative derivation of Manyama's (2026) convolution,
complementing Andrews--El Bachraoui's (2026) strict-partition formulas.
An injection sends a partition with nonnegative
reverse-alternating sum to an odd-length partition of the next weight.
We characterize equality, give an explicit lower bound for the deficit,
and deduce monotonicity of the corresponding unrestricted count.
The same parity analysis proves a shifted companion comparison, determines
all plateaus of the negative-sum count, and corrects neighboring indexing
and distinct-part descriptions. A bivariate series settles the linked
ordered-composition interpretations, with their precise shifts.
\end{abstract}

\section{The question and the partition counts}
Wiseman's FORMULA in OEIS A344611, revision 27, states
``Conjecture: a(n) <= A160786(n)'' \cite{a344611}.
Here the first count concerns partitions of $2n$ with nonnegative
reverse-alternating sum, and the second counts partitions of $2n+1$
with an odd number of parts \cite{a160786,a027193}.
We prove the comparison, including $n=0$, by an explicit injection.
Its missing images also determine exactly when equality holds.

A partition $\lambda=(\lambda_1,\ldots,\lambda_k)$ has positive integer
parts in nonincreasing order, weight $|\lambda|=\sum_i\lambda_i$,
and length $\len(\lambda)=k$. Put
\[
 R(\lambda)=\sum_{i=1}^k(-1)^{k-i}\lambda_i,
 \qquad R(\varnothing)=0.
\]
Let $O(m)$ and $E(m)$ count partitions of $m$ of odd and even length,
respectively, and let $p(m)=O(m)+E(m)$ \cite{a027193,a027187,a000041}.
The empty partition has even
length, so $p(0)=E(0)=1$ and $O(0)=0$. Write
\[
 C(m)=\#\{\lambda:|\lambda|=m,\ R(\lambda)\ge0\},\qquad
 N(m)=\#\{\lambda:|\lambda|=m,\ R(\lambda)<0\}.
\]
Thus A344611$(n)=C(2n)$, A160786$(n)=O(2n+1)$,
A344743$(n)=N(2n)$, and A236914$(n)=E(2n+1)$
\cite{a344611,a160786,a344743,a236914}.

The distinct-part companion A344650 occurs in the convolution below.
Andrews and El Bachraoui (2026), Section 3, Theorem 5, give generating
functions for this distinct-part count \cite{andrews}.
Their statement is different from the unrestricted-count inequality
proved here. The convolution recorded by Seiichi Manyama on May 17, 2026,
and the conjugation interpretation listed in A344611 receive
elementary proofs below. These identities supply the distinct-part
setting for the injection and its signed-endpoint extension.

\section{Parity and the injection}
\begin{lemma}\label{lem:sign}
Every odd-length partition has $R(\lambda)>0$.
An even-length partition has $R(\lambda)\le0$, with equality precisely
when $\lambda_{2i-1}=\lambda_{2i}$ for every adjacent pair.
More precisely, if $o(\lambda')$ is the number of odd parts of the
conjugate, then $R(\lambda)=(-1)^{k-1}o(\lambda')$ for $k>0$.
Consequently,
\begin{equation}\label{eq:counts}
 C(m)=O(m)+\begin{cases}p(m/2),&m\text{ even},\\0,&m\text{ odd},\end{cases}
 \quad
 N(m)=E(m)-\begin{cases}p(m/2),&m\text{ even},\\0,&m\text{ odd}.
 \end{cases}
\end{equation}
\end{lemma}
\begin{proof}
For $k$ odd,
$R(\lambda)=(\lambda_1-\lambda_2)+\cdots+
(\lambda_{k-2}-\lambda_{k-1})+\lambda_k>0$.
For $k$ even,
$R(\lambda)=-\sum_{i=1}^{k/2}(\lambda_{2i-1}-\lambda_{2i})$.
Each difference is nonnegative, so equality means that all these
differences vanish. The number of odd-height columns is
$\sum_i(\lambda_{2i-1}-\lambda_{2i})$, with an unpaired last part
included when $k$ is odd; this proves the conjugate formula.
Such zero-sum partitions are exactly the partitions obtained
by duplicating every part of an arbitrary partition of $m/2$.
The empty partition is included in this description.
\end{proof}

\begin{theorem}\label{thm:main}
For every $m\in\NN$, $C(m)\le O(m+1)$.
In particular, A344611$(n)\le$ A160786$(n)$ for every $n\in\NN$.
Equality in this latter comparison holds exactly for $0\le n\le3$.
For $n\ge4$, the deficit satisfies
\begin{equation}\label{eq:gap}
 O(2n+1)-C(2n)\ \ge\ \left\lfloor\frac{n-4}{2}\right\rfloor+1.
\end{equation}
\end{theorem}
\begin{proof}
For a partition counted by $C(m)$ define
\[
 F(\lambda)=
 \begin{cases}
 (\lambda_1+1,\lambda_2,\ldots,\lambda_k),&k\text{ odd},\\
 (\lambda_1,\ldots,\lambda_k,1),&k\text{ even}.
 \end{cases}
\]
In particular $F(\varnothing)=(1)$. Both branches preserve the
nonincreasing order, increase weight by one, and have odd length.
Each branch is injective: subtract one from the first part in the
odd branch, or delete the last $1$ in the even branch.

The odd branch has a unique largest part and that part is at least $2$.
By Lemma~\ref{lem:sign}, every nonempty even source starts with two equal
parts, so its image has a repeated largest part. The exceptional even
source $\varnothing$ has image $(1)$, which cannot come from the odd
branch. Thus the images of the branches are disjoint, proving injectivity.

For $n\ge4$ and $0\le j\le\lfloor(n-4)/2\rfloor$, consider
\[
 \mu^{(j)}=(3,3,3,\underbrace{2,\ldots,2}_{2j},
                   \underbrace{1,\ldots,1}_{2n-8-4j}).
\]
These are distinct odd-length partitions of $2n+1$ with repeated largest
part and $R(\mu^{(j)})=3$. They cannot lie in the odd branch. Every even
branch image has reverse-alternating sum $1$, since its source has sum
$R=0$ and appending $1$ changes $R$ to $1-R$. Hence none of the
$\mu^{(j)}$ lies in the image of $F$, proving \eqref{eq:gap}.
For $n=0,1,2,3$, the two counts are respectively $1,2,4,8$,
as direct enumeration of partitions of weights at most $7$ gives.
\end{proof}

For completeness, equality in the general-weight comparison is exactly
$m\in\{0,1,2,3,4,6\}$. The even weights were just treated. At odd
$m\ge5$, the target $(2,2,2,1^{m-5})$ has odd length and repeated
largest part, so is outside the odd branch; an even paired source
cannot have odd weight. At $m=1,3$ the respective counts are $1,2$.

\begin{corollary}\label{cor:mono}
The sequence $C(m)$, namely A344607, is nondecreasing,
answering the question in its COMMENTS section \cite{a344607}.
Its adjacent equalities are exactly
$C(m)=C(m+1)$ for $m\in\{0,2,4,6\}$.
\end{corollary}
\begin{proof}
For even $m$, \eqref{eq:counts} gives $C(m+1)=O(m+1)$,
so the conclusion is Theorem~\ref{thm:main}.
For odd $m$, the same formula and that theorem give
$C(m)\le O(m+1)<O(m+1)+p((m+1)/2)=C(m+1)$.
Here $p((m+1)/2)\ge1$.
\end{proof}

\section{The companion comparison and two source identities}
\begin{proposition}\label{prop:negative}
The negative-sum counts, namely A344608, satisfy
$N(m)\le N(m+1)$ for every $m\ge0$, answering its COMMENTS
monotonicity question \cite{a344608}.
The COMMENTS of A344743 say exactly ``Conjecture: a(n) >= A236914.''
\cite{a344743}, without specifying the argument on the right.
A common-index reading is not asserted explicitly in the source;
it fails at $n=1$, with sides $0$ and $1$.
A precise valid comparison is
\begin{equation}\label{eq:corrected}
 \mathrm{A344743}(n+1)\ \ge\ \mathrm{A236914}(n)
 \qquad(n\ge0).
\end{equation}
Equality in \eqref{eq:corrected} holds exactly for $n=0,1,2,3$;
for $n\ge4$ its deficit is at least $\lfloor(n-4)/2\rfloor+1$.
More generally, $N(m)=N(m+1)$ exactly for
$m\in\{0,1,3,5,7\}$.
By conjugation, $N(m)$ also counts partitions with even largest part
and at least one odd part. In particular, the odd-weight identity is
$N(2n+1)=E(2n+1)$; the index $2n+2$ and the phrase concerning parts
not all odd in the COMMENTS of A344608 require these corrections.
\end{proposition}
\begin{proof}
By Lemma~\ref{lem:sign}, a negative-sum partition has even length.
Increasing its first part by $1$ changes $R$ to $R-1<0$ and gives
an injective map to negative-sum partitions of the next weight.
This proves monotonicity.
Every even-length partition of odd weight has $R<0$, since a partition
with all adjacent pairs equal has even weight. Thus $N(2n+1)=E(2n+1)$,
and monotonicity proves \eqref{eq:corrected}.
At $n=1$, the partitions of $2$ are $(2)$ and $(1,1)$, with sums $2$
and $0$, so $N(2)=0$; the only even-length partition of $3$ is $(2,1)$,
so $E(3)=1$. Thus the common-index reading fails.
For $n\ge4$ and $0\le j\le\lfloor(n-4)/2\rfloor$,
$(3,3,3,2^{2j},1^{2n-7-4j})$ is a partition of $2n+2$
of even length with $R=-2$. These distinct targets have repeated
largest part, whereas the increase-first-part map has a unique largest
part. They prove the quantitative bound. For even $m=2r\ge2$,
the target $(r+1,r)$ has $R=-1$; its only possible preimage $(r,r)$
has sum zero. Hence $N(m)<N(m+1)$. For odd $m\ge9$ the preceding
witness applies. Direct enumeration gives the remaining values
$N(0),\ldots,N(8)=0,0,0,1,1,3,3,7,7$, establishing both equality claims.
For the conjugation assertion, even length is equivalent to even
largest part of the conjugate. Within this class $R<0$ means that
some adjacent pair of original parts is unequal, or equivalently
that a column has odd height. This means exactly that the conjugate
has at least one odd part.
\end{proof}

\begin{corollary}\label{cor:distinct}
Among distinct-part partitions, $R>0$ means odd length and $R<0$
means even length. Thus the distinct negative-sum count is
\textup{A067661}$(m)$ for $m>0$, and is zero at $m=0$.
The distinct nonnegative-sum count is \textup{A067659}$(m)$ for $m>0$,
with the empty partition added at zero.
The identification of the distinct negative-sum count with A067659
in A344608, CROSSREFS, fails at $m=5$: $(4,1),(3,2)$ give two
negative objects, whereas $(5)$ is the sole positive one.
\end{corollary}
\begin{proof}
Every difference in Lemma~\ref{lem:sign} is positive for distinct parts.
Writing $D_+(x)=\prod_{s\ge1}(1+x^s)$ and
$D_-(x)=\prod_{s\ge1}(1-x^s)$, the even- and odd-length generating
functions are $(D_++D_-)/2$ and $(D_+-D_-)/2$.
This rederives the formulas recorded by Jovovic (2002)
\cite{a067659,a067661}; the latter also occurs in
Andrews--El Bachraoui (2026), Theorem 5 \cite{andrews}.
\end{proof}

\begin{proposition}\label{prop:identities}
Let $d(0)=1$, and for $r\ge1$ let $d(r)$ count distinct-part partitions
of $2r$ with odd length, namely A344650$(r)$.
Then
\begin{equation}\label{eq:convolution}
 C(2n)=\sum_{r=0}^n d(r)p(n-r).
\end{equation}
Conjugation also identifies the partitions counted by $C(m)$ with
partitions of $m$ whose largest part is odd or whose parts are all even,
with the empty partition included in the latter class.
More generally, let $B_m(s)$ and $D_m(s)$ count unrestricted and
distinct-part partitions of $m$ with $R=s$. Then for every integer $s$,
\begin{equation}\label{eq:signed}
 B_m(s)=\sum_{t=0}^{\lfloor m/2\rfloor}p(t)D_{m-2t}(s).
\end{equation}
\end{proposition}
\begin{proof}
For each part size $s$, write its multiplicity uniquely as
$\epsilon_s+2t_s$ with $\epsilon_s\in\{0,1\}$.
This decomposes a partition as a distinct-part partition $\rho$ and
a partition $\nu$ with every part duplicated. Removing two equal
adjacent parts cancels their contributions to $R$, and preserves the
signs of all remaining parts. Thus $R(\lambda)=R(\rho)$ and
$|\lambda|=|\rho|+2|\nu|$, proving \eqref{eq:signed} bijectively.
The parity of the original
length equals that of $\len(\rho)$. Lemma~\ref{lem:sign} therefore says
that a partition of $2n$ is counted by $C(2n)$ exactly when
$\rho$ is empty or has odd length. In either case $|\rho|=2r$ is even,
and $|\nu|=n-r$. This is a bijective proof of \eqref{eq:convolution},
the convolution displayed in A344611 \cite{a344611,a344650}.
The value $d(0)=1$ represents the empty residual partition, not an
odd-length partition of zero.

For conjugation, the largest part of $\lambda'$ is $\len(\lambda)$.
An odd-length $\lambda$ therefore corresponds exactly to a conjugate
with odd largest part. A partition with all adjacent pairs equal
corresponds exactly to a conjugate whose parts are all even: the latter
parts are the column heights of the original Ferrers diagram.
These two classes are disjoint. Lemma~\ref{lem:sign} gives the result.
\end{proof}

Here $C(2n)$ is the \emph{even} bisection of A344607; the word
``Odd'' in A344611, CROSSREFS, needs this replacement.
Also $O(2n)=\mathrm{A236559}(n)$, so
$C(2n)=\mathrm{A236559}(n)+p(n)$, rather than an equality of those
two sequences \cite{a236559}.
Equation~\eqref{eq:signed} rederives Manyama's (May 17, 2026)
column convolution in A344610 from A344649, and extends it to the
signed triangles A344612 and A344739 \cite{triangles}.
For the even-minus-odd indexed-part expression in A344612, COMMENTS,
the factor is $(-1)^k$, not $(-1)^{k-1}$: the singleton $(1)$ has
$R=1$, not $-1$.

\begin{corollary}\label{cor:ordered}
Let $H(m)$ and $J(m)$ count compositions of $m$ with reverse-alternating
sum nonnegative and negative, respectively. Then $H(0)=1$, $J(0)=0$,
and for $m\ge1$,
\[
 H(m)=\sum_{j=0}^{\lfloor m/2\rfloor}\binom{m-1}{j}
 =\mathrm{A116406}(m),\qquad
 J(m)=2^{m-1}-H(m)=\mathrm{A294175}(m-1).
\]
In particular $H(2r)=\mathrm{A114121}(r)$ for $r\ge0$, and
$J(2r)=\mathrm{A008549}(r-1)$ for $r\ge1$.
These prove the ``appears to be'' ordered interpretations in
A344607 and A344611, and specify the shift missing from A344608.
\end{corollary}
\begin{proof}
Reversing a composition changes its reverse-alternating sum into the
ordinary alternating sum. Mark weight by $x$ and this latter sum by $z$.
With $f(t)=t/(1-t)$, grouping consecutive positive and negative positions
gives the formal series
\[
 \frac{1+f(xz)}{1-f(xz)f(x/z)}
 =\frac{1-x/z}{1-x(z+z^{-1})}.
\]
For $m\ge1$, its $x^m z^s$ coefficient is
$\binom{m-1}{(m+s)/2-1}$, with an out-of-range or noninteger
lower argument interpreted as zero. Summing over $s\ge0$ and using
binomial symmetry gives $H(m)$ above. The complementary count is $J(m)$.
The formulas in A116406 (Barry, 2007), A114121 (Barry, 2009),
A294175 (Israel, 2018), and A008549 give the stated identifications
\cite{a116406,a114121,a294175,a008549}. For example,
$H(2r)=4^{r-1}+\frac12\binom{2r}{r}$ and
$J(2r)=4^{r-1}-\binom{2r-1}{r-1}$ for $r\ge1$.
The unshifted formula $H(m)=2^{m-1}-\mathrm{A294175}(m)$ in
A116406, FORMULA, is corrected by $m-1$: at $m=4$ it is $3$ while
$H(4)=7$. Likewise, the complementary formula in A114121, FORMULA
(Wiseman, June 27, 2021), needs A008549$(r-1)$ instead of A008549$(r)$;
at $r=2$ the unshifted right side is $2$ rather than $7$.
\end{proof}

The formula for $H(m)$ also counts binary words of length $m$ beginning
with $1$ and having at most $\lfloor m/2\rfloor$ zeros. In the path
description in A116406, COMMENTS (Stephan, 2008), the corresponding
condition is a nonnegative \emph{endpoint}, not nonnegative height
at every step: $1001$ contributes at length four but its path visits
height $-1$. There are seven such words and only six nonnegative paths
of length four.

\begin{corollary}\label{cor:recurrence}
The recurrence conjectured by Mathar (2014) in A116406, FORMULA, holds
for $m\ge3$:
\[
 m(m-3)H(m)+2(-m^2+4m-2)H(m-1)
 -4(m-2)^2H(m-2)+8(m-2)(m-3)H(m-3)=0.
\]
\end{corollary}
\begin{proof}
The binomial formulas give
$G(x)=\sum_{m\ge0}H(m)x^m
=(1+x-2x^2)/(2(1-4x^2))+(1+x)/(2\sqrt{1-4x^2})$.
With $\theta=x\,d/dx$, differentiation gives
\[
 [(1-2x-4x^2+8x^3)\theta^2+(-3+4x+8x^3)\theta+2x]G=0.
\]
The coefficient of $x^m$ is exactly the stated recurrence.
\end{proof}

\section{Coverage and exact checks}
\begin{center}\footnotesize
\begin{tabular}{p{0.37\linewidth}p{0.54\linewidth}}
\toprule
Source item & Coverage \\
\midrule
A344611, FORMULA comparison & Proved; Theorem~\ref{thm:main}. \\
A344611, displayed deficit & Not posed as an all-index formula;
equality exactly for $n=0,1,2,3$;
strictness and bound \eqref{eq:gap} for $n\ge4$. \\
A344611, convolution & Reproved; prior formula Manyama (2026);
Proposition~\ref{prop:identities}. \\
A344610 column convolution; A344612/A344739 signed triangles &
Reproved; prior formula Manyama (2026) for A344610;
all weights and signed endpoints covered by \eqref{eq:signed}. \\
A344612, indexed-part expression & Refuted-and-corrected at $(1)$;
factor $(-1)^k$ for even-minus-odd indexed sums. \\
A344611/A344607, conjugation comment & Proved; Lemma~\ref{lem:sign}
and Proposition~\ref{prop:identities}. \\
A344607 monotonicity & Proved with exact plateaus;
Corollary~\ref{cor:mono}. \\
A344743, COMMENTS comparison & Argument unspecified; common-index
reading not asserted explicitly, fails at $n=1$; shifted statement
\eqref{eq:corrected} proved (not posed in that form). \\
A344608 monotonicity, odd-weight bisection, conjugation & Monotonicity
proved with exact plateaus; index and parity wording refuted-and-corrected;
Proposition~\ref{prop:negative}. \\
A344608, distinct-part cross-reference & Refuted-and-corrected;
Corollary~\ref{cor:distinct}. \\
A344611, bisection cross-reference & Refuted-and-corrected: even bisection. \\
A344650 generating functions & Prior formulas cited:
Andrews--El Bachraoui (2026), Theorem 5, (3.5). Its odd-length
product (3.4) is reproved in Corollary~\ref{cor:distinct};
prior formula Jovovic (2002). \\
A344607/A344611/A344608 ordered interpretations & Proved with exact shifts;
Corollary~\ref{cor:ordered}; formulas rederived from Barry (2007, 2009)
and Israel (2018). \\
A116406/A114121, complementary formulas & Refuted-and-corrected;
Corollary~\ref{cor:ordered}. \\
A116406, path equivalence and recurrence & Path condition
refuted-and-corrected at length four; Mathar's (2014) recurrence proved
in Corollary~\ref{cor:recurrence}. \\
A114121, recurrence & Prior proof cited: Israel (2020), FORMULA. \\
A236559, character-table conjecture;
Andrews--El Bachraoui (2026), Section 11; further asymptotics and
column formulas in A344610/A344649 & Not covered: different objects
or additional identities. \\
\bottomrule
\end{tabular}
\end{center}

The declaration \texttt{result} in the pinned Lean source \cite{lean}
has conclusion $\forall n\in\NN,\ C(2n)\le O(2n+1)$.
Its definitions use positive multiset parts sorted in decreasing order
and the endpoint recurrence $s_0=0$, $s_i=\lambda_i-s_{i-1}$,
whose last value is $R(\lambda)$. Thus its source and target counts
match those above. Its dependency \texttt{trimmedSums} in
\href{https://github.com/the-omega-institute/trureturing/blob/d133fcb1d5e4878bdf5232fd0c07ef2abb3da1d9/D5/S1/Words/Compositions/TrimmedAlternatingPartitions.lean}{\texttt{TrimmedAlternatingPartitions.lean}}
is the recurrence $s_0=0$, $s_i=\lambda_i-s_{i-1}$; induction gives
$s_k=R(\lambda)$. The declaration
\texttt{trimmedSums\_nodup\_iff\_strict\_tail} states that for positive
nonincreasing parts, the recurrence endpoints are distinct exactly
when the tail is strictly decreasing. Its
\texttt{card\_trimmedSums\_eq\_distincts} equates that endpoint-distinct
count at weight $m$ with distinct-part partitions of weight $m+1$;
neither is the sign comparison.
The general-weight theorem, strictness, plateaus, companion corrections,
source identities, and ordered compositions have written proofs here.

The accompanying \texttt{verify.py} enumerates all partitions of weights
$0$ through $42$, checks both maps through source weight $41$, and
compares the counts with independent length-parity dynamic programs.
It checks \eqref{eq:counts}, \eqref{eq:gap}, and
\eqref{eq:convolution}, the shifted comparison, the conjugation
interpretation, and the displayed OEIS initial values through $n=20$.
It also checks the exact plateau sets, distinct-part corrections,
the signed convolution, and all compositions of weights $0$ through
$16$ against the bivariate coefficient formula and the cited initial
values. The ordered recurrence is checked through index $100$.
These are finite checks; the unbounded assertions follow from the
proofs above.

\begin{thebibliography}{20}
\footnotesize
\setlength{\itemsep}{0pt}
\bibitem{a344611} G. Wiseman, OEIS A344611 (2021),
FORMULA, COMMENTS, EXAMPLE, and AUTHOR,
\url{https://oeis.org/A344611}.
\bibitem{a160786} U. Sarkar, OEIS A160786 (2009), NAME and FORMULA,
\url{https://oeis.org/A160786}.
\bibitem{a027193} OEIS Foundation Inc., OEIS A027193, NAME,
\url{https://oeis.org/A027193}.
\bibitem{a027187} OEIS Foundation Inc., OEIS A027187, NAME,
\url{https://oeis.org/A027187}.
\bibitem{a000041} OEIS Foundation Inc., OEIS A000041, NAME,
\url{https://oeis.org/A000041}.
\bibitem{a344743} G. Wiseman, OEIS A344743 (2021), COMMENTS,
\url{https://oeis.org/A344743}.
\bibitem{a236914} C. Kimberling, OEIS A236914 (2014), NAME and COMMENTS,
\url{https://oeis.org/A236914}.
\bibitem{a344650} G. Wiseman, OEIS A344650 (2021), NAME and FORMULA,
\url{https://oeis.org/A344650}.
\bibitem{andrews} G. E. Andrews and M. El Bachraoui,
\emph{On a two-color partition series and its companions},
arXiv:2606.30208v1 (2026), Section 3, Theorem 5, pp.~6 and 16--17,
\url{https://arxiv.org/abs/2606.30208v1}.
Journal metadata: \emph{J. Number Theory} \textbf{292} (2027), 78--102,
\href{https://doi.org/10.1016/j.jnt.2026.06.002}{doi:10.1016/j.jnt.2026.06.002};
the theorem locators above refer to the preprint.
\bibitem{lean} \emph{WisemanA344611OddLengthInjection.lean},
declaration \texttt{result}, commit \texttt{d133fcb1d5e4},
\href{https://github.com/the-omega-institute/trureturing/blob/d133fcb1d5e4878bdf5232fd0c07ef2abb3da1d9/D5/S1/Words/Compositions/WisemanA344611OddLengthInjection.lean}{pinned Lean declaration}.
\bibitem{a344607} G. Wiseman, OEIS A344607 (2021), COMMENTS,
\url{https://oeis.org/A344607}.
\bibitem{a344608} G. Wiseman, OEIS A344608 (2021), COMMENTS,
\url{https://oeis.org/A344608}.
\bibitem{a067659} OEIS A067659, NAME and FORMULA (Jovovic, February 24, 2002),
\url{https://oeis.org/A067659}.
\bibitem{a067661} OEIS A067661, NAME and FORMULA (Jovovic, February 24, 2002),
\url{https://oeis.org/A067661}.
\bibitem{a236559} C. Kimberling, OEIS A236559 (2014), COMMENTS,
\url{https://oeis.org/A236559}.
\bibitem{a116406} P. Barry, OEIS A116406, FORMULA (October 6, 2007),
and G. Wiseman, COMMENTS (June 20, 2021), \url{https://oeis.org/A116406}.
\bibitem{a114121} OEIS A114121, FORMULA (Barry, September 7, 2009;
Israel, July 27, 2020),
COMMENTS (Wiseman, June 21, 2021), \url{https://oeis.org/A114121}.
\bibitem{a294175} OEIS A294175, FORMULA (Israel, February 12, 2018),
COMMENTS (Wiseman, July 22, 2021), \url{https://oeis.org/A294175}.
\bibitem{a008549} OEIS A008549, FORMULA and OFFSET,
\url{https://oeis.org/A008549}.
\bibitem{triangles} G. Wiseman, OEIS signed-partition triangles (2021):
\href{https://oeis.org/A344610}{A344610} (FORMULA, Manyama, May 17, 2026),
\href{https://oeis.org/A344649}{A344649},
\href{https://oeis.org/A344612}{A344612}, and
\href{https://oeis.org/A344739}{A344739} (NAME and COMMENTS).
\end{thebibliography}
\noindent\footnotesize Sources accessed November 4--5, 2026. License:
\href{https://creativecommons.org/licenses/by/4.0/}{Creative Commons Attribution 4.0 (CC BY 4.0)}.
\end{document}
